\section{Подготовка: где писать и запускать код}
\label{sec:start}

\begin{infobox}[Цели главы]
\begin{itemize}
    \item Понять, чем отличаются запуск в симуляторе и на настоящем дроне, и почему код для них один.
    \item Найти в симуляторе редактор, выбор языка, кнопки запуска и терминал.
    \item Запустить первую программу и увидеть её вывод.
    \item Запомнить правила безопасности при работе с настоящим Пионером.
\end{itemize}
\end{infobox}

\subsection{Два способа запустить программу}

У программы для Пионера есть два «исполнителя»:

\begin{center}
\small
\begin{tabularx}{\linewidth}{@{}>{\bfseries}l X X@{}}
\toprule
 & Симулятор в браузере & Настоящий Пионер \\
\midrule
Что нужно & Любой компьютер с браузером: \url{https://simulator.sakoryagin.ru} & Дрон, аккумулятор, кабель USB или Wi-Fi, площадка \\
Lua & Редактор на сайте, язык \ui{Lua} & Pioneer Station: загрузка скрипта на борт \\
Python & Редактор на сайте, язык \ui{Python} & \texttt{pip install pioneer-sdk}, программа на вашем компьютере \\
Что видно & Полёт в 3D, терминал, телеметрия & Сам дрон и консоль Pioneer Station \\
Цена ошибки & Нажать \ui{Сбросить} & Сломанный пропеллер или хуже \\
\bottomrule
\end{tabularx}
\end{center}

Код в обоих случаях один и тот же: симулятор воспроизводит официальный API автопилота (\texttt{ap}, \texttt{Ev}, \texttt{Timer}, \texttt{Ledbar}, \texttt{Sensors}) и библиотеку \texttt{pioneer\_sdk}. Поэтому в курсе мы сначала отлаживаем программу в симуляторе, а на дрон переносим уже проверенную.

\begin{rulebox}[Сначала симулятор]
Любая новая программа запускается сначала в симуляторе. На настоящем дроне запускается только программа, которая в симуляторе взлетела, выполнила задание и села.
\end{rulebox}
